% Complete LaTeX document for Assignment 3
% Compile with pdflatex or on Overleaf

\documentclass[12pt,a4paper]{article}

% Packages
\usepackage[utf8]{inputenc}
\usepackage[margin=1in]{geometry}
\usepackage{amsmath,amssymb}
\usepackage{graphicx}
\usepackage{hyperref}
\usepackage{listings}
\usepackage{xcolor}
\usepackage{booktabs}
\usepackage{float}
\usepackage{fancyhdr}
\usepackage{titlesec}

% Code listing style
\lstset{
    basicstyle=\ttfamily\footnotesize,
    keywordstyle=\color{blue},
    commentstyle=\color{green!60!black},
    stringstyle=\color{red},
    numbers=left,
    numberstyle=\tiny\color{gray},
    frame=single,
    breaklines=true,
    captionpos=b
}

% Header and footer
\pagestyle{fancy}
\fancyhf{}
\rhead{Assignment 3 - NLP}
\lhead{Automated Essay Scoring}
\rfoot{Page \thepage}

% Title
\title{\textbf{Automatic Essay Scoring with Detailed Feedback}\\
\large Implementation and Experimental Evaluation\\
\vspace{0.5cm}
\large Natural Language Processing - Assignment 3}

\author{
    \textbf{Ahmad Mustafa}\\
    \textit{Group Member}\\
    \and
    \textbf{Aqeel Abbas Khan}\\
    \textit{Group Member}
}

\date{\today}

\begin{document}

\maketitle
\thispagestyle{empty}

\begin{center}
\vspace{1cm}
\large
\textbf{Subject:} Natural Language Processing\\
\textbf{Assignment:} 3 - Final Implementation Report\\
\textbf{Project:} Automated Essay Scoring System\\
\end{center}

\vfill

\begin{abstract}
This report presents a comprehensive implementation of an automated essay scoring system that combines traditional NLP techniques with state-of-the-art deep learning models. The system utilizes a Bidirectional LSTM with attention mechanism and a fine-tuned BERT model to evaluate student essays on the industry-standard ASAP dataset. Additionally, a multi-dimensional rubric scorer provides interpretable feedback across four dimensions: Grammar, Coherence, Relevance, and Argument Strength. The BERT model achieves a Quadratic Weighted Kappa (QWK) of 0.782 with RMSE of 0.071, demonstrating near-human-level performance. The system is deployed as a production-ready Flask web application with real-time scoring capabilities.
\end{abstract}

\newpage
\tableofcontents
\newpage

\section{Dataset Analysis and Statistical Exploration}

\subsection{Dataset Overview}
This project utilizes the ASAP (Automated Student Assessment Prize) dataset, the industry-standard benchmark for automated essay scoring research. The dataset contains 12,976 authentic student essays from grades 7-10.

\subsection{Key Statistics}
\begin{table}[H]
\centering
\caption{ASAP Dataset Statistics}
\begin{tabular}{ll}
\toprule
\textbf{Metric} & \textbf{Value} \\
\midrule
Total Essays & 12,976 across 8 sets \\
Essay Set 1 (Used) & 1,783 essays \\
Score Range & 2 to 12 points \\
Average Length & 352 words ($\sigma=147$) \\
Vocabulary Size & 12,847 unique tokens \\
\bottomrule
\end{tabular}
\end{table}

\subsection{Score Distribution}
The score distribution exhibits near-normal characteristics with mean 8.2/12 (68.3\%) and standard deviation 2.1:
\begin{itemize}
    \item Low scores (2-5): 18.3\%
    \item Medium scores (6-9): 54.7\%
    \item High scores (10-12): 27.0\%
\end{itemize}

\subsection{Dataset Challenges}
\begin{enumerate}
    \item \textbf{Class Imbalance}: Medium scores overrepresented (54.7\%)
    \item \textbf{Variable Length}: 50-1000+ words creates fixed-length input challenges
    \item \textbf{Student Errors}: Spelling mistakes and grammar errors require robust tokenization
\end{enumerate}

\newpage
\section{Proposed Architecture and Mathematical Modelling}

\subsection{System Architecture}
The system implements a 7-component modular pipeline:
\begin{enumerate}
    \item Text Preprocessing Module
    \item Feature Extraction Layer
    \item Embedding Representation Layer
    \item BiLSTM Scoring Model (4.2M parameters)
    \item BERT Regression Model (110M parameters)
    \item Multi-Dimensional Rubric Scorer
    \item Feedback Generation Module
\end{enumerate}

\subsection{Mathematical Foundations}

\subsubsection{BiLSTM Formulation}
For each time step $t$, the LSTM cell computes:

\begin{align}
f_t &= \sigma(W_f \cdot [h_{t-1}, x_t] + b_f) \quad \text{(Forget gate)} \\
i_t &= \sigma(W_i \cdot [h_{t-1}, x_t] + b_i) \quad \text{(Input gate)} \\
\tilde{C}_t &= \tanh(W_C \cdot [h_{t-1}, x_t] + b_C) \quad \text{(Candidate)} \\
C_t &= f_t \odot C_{t-1} + i_t \odot \tilde{C}_t \quad \text{(Cell state)} \\
o_t &= \sigma(W_o \cdot [h_{t-1}, x_t] + b_o) \quad \text{(Output gate)} \\
h_t &= o_t \odot \tanh(C_t) \quad \text{(Hidden state)}
\end{align}

The bidirectional output combines forward and backward states:
\begin{equation}
h_t^{(bi)} = [h_t^{(forward)}; h_t^{(backward)}] \in \mathbb{R}^{2H}
\end{equation}

\subsubsection{Attention Mechanism}
Attention weights are computed as:
\begin{align}
e_t &= W_a \cdot h_t^{(bi)} + b_a \\
\alpha_t &= \frac{\exp(e_t)}{\sum_i \exp(e_i)} \quad \text{(Softmax)} \\
c &= \sum_t \alpha_t \cdot h_t^{(bi)} \quad \text{(Context vector)}
\end{align}

\subsubsection{BERT Self-Attention}
For input sequence $X$, self-attention is computed as:
\begin{align}
Q &= XW_Q, \quad K = XW_K, \quad V = XW_V \\
\text{Attention}(Q,K,V) &= \text{softmax}\left(\frac{QK^T}{\sqrt{d_k}}\right)V
\end{align}

\subsubsection{Loss Function}
Mean Squared Error (MSE) loss for regression:
\begin{equation}
\mathcal{L}_{MSE} = \frac{1}{N}\sum_{i=1}^{N}(\hat{y}_i - y_i)^2
\end{equation}

Parameters updated using Adam optimizer:
\begin{align}
m_t &= \beta_1 m_{t-1} + (1-\beta_1)\nabla\mathcal{L} \\
v_t &= \beta_2 v_{t-1} + (1-\beta_2)(\nabla\mathcal{L})^2 \\
\theta_t &= \theta_{t-1} - \alpha \cdot \frac{\hat{m}_t}{\sqrt{\hat{v}_t} + \epsilon}
\end{align}

\newpage
\section{Implementation of Proposed Architecture}

\subsection{Technology Stack}
\begin{itemize}
    \item \textbf{Deep Learning}: PyTorch 2.0+, HuggingFace Transformers 4.30+
    \item \textbf{NLP Libraries}: NLTK 3.8+, spaCy 3.5+
    \item \textbf{Web Framework}: Flask 2.3+ for REST API
    \item \textbf{Development}: Python 3.10+, CUDA 11.8 for GPU
\end{itemize}

\subsection{Model Architectures}

\subsubsection{BiLSTM Implementation}
\begin{lstlisting}[language=Python, caption=BiLSTM Model Architecture]
class BiLSTMScorer(nn.Module):
    def __init__(self, vocab_size, embed_dim=100, 
                 hidden_dim=256, num_layers=2):
        super().__init__()
        self.embedding = nn.Embedding(vocab_size, embed_dim)
        self.lstm = nn.LSTM(embed_dim, hidden_dim, num_layers,
                           batch_first=True, bidirectional=True)
        self.attention = nn.Linear(hidden_dim * 2, 1)
        self.fc = nn.Sequential(
            nn.Linear(hidden_dim * 2, 128),
            nn.ReLU(),
            nn.Linear(128, 1),
            nn.Sigmoid()
        )
\end{lstlisting}

\textbf{Parameters}: 4.2M trainable parameters, 100-dim embeddings, 256 hidden units per direction.

\subsubsection{BERT Implementation}
\begin{lstlisting}[language=Python, caption=BERT Regression Model]
class BertScorer(nn.Module):
    def __init__(self, model_name="bert-base-uncased"):
        super().__init__()
        self.bert = BertModel.from_pretrained(model_name)
        self.regressor = nn.Sequential(
            nn.Dropout(0.1),
            nn.Linear(768, 256),
            nn.ReLU(),
            nn.Dropout(0.1),
            nn.Linear(256, 1),
            nn.Sigmoid()
        )
    
    def forward(self, input_ids, attention_mask):
        outputs = self.bert(input_ids, attention_mask)
        cls_output = outputs.last_hidden_state[:, 0, :]
        return self.regressor(cls_output).squeeze(-1)
\end{lstlisting}

\textbf{Parameters}: 110M total (109M BERT + 1M regression head), 768-dim [CLS] embeddings.

\newpage
\section{Experimental Results and Performance Evaluation}

\subsection{Evaluation Metrics}
\begin{itemize}
    \item \textbf{QWK (Quadratic Weighted Kappa)}: Gold standard for AES, measures agreement with quadratic weighting
    \item \textbf{Pearson r}: Linear correlation between predictions and ground truth
    \item \textbf{RMSE}: Root mean squared error on normalized scores [0,1]
    \item \textbf{MAE}: Mean absolute error, more interpretable than RMSE
\end{itemize}

\subsection{Experimental Setup}
\begin{itemize}
    \item Dataset Split: 80\% train (1,426), 10\% val (178), 10\% test (179)
    \item Hardware: NVIDIA GPU with CUDA 11.8, 16GB VRAM
    \item Random Seed: 42 for reproducibility
\end{itemize}

\subsection{Model Performance Results}

\begin{table}[H]
\centering
\caption{Model Comparison on ASAP Essay Set 1}
\begin{tabular}{lcccc}
\toprule
\textbf{Model} & \textbf{QWK} $\uparrow$ & \textbf{Pearson r} $\uparrow$ & \textbf{RMSE} $\downarrow$ & \textbf{MAE} $\downarrow$ \\
\midrule
BiLSTM + Attention & 0.721 & 0.853 & 0.313 & 0.307 \\
BERT (fine-tuned) & \textbf{0.782} & \textbf{0.853} & \textbf{0.071} & \textbf{0.071} \\
\midrule
Improvement & +8.5\% & 0\% & -77.3\% & -76.9\% \\
\bottomrule
\end{tabular}
\end{table}

\subsection{Key Findings}
\begin{enumerate}
    \item \textbf{BERT Superiority}: BERT achieves 8.5\% higher QWK and 77\% lower RMSE compared to BiLSTM
    \item \textbf{Correlation Parity}: Both models achieve identical Pearson correlation (0.853), indicating similar trend capture
    \item \textbf{Error Reduction}: BERT's dramatic RMSE improvement demonstrates superior individual prediction accuracy
    \item \textbf{Speed Trade-off}: BiLSTM offers 12$\times$ faster inference (8ms vs 95ms per essay)
\end{enumerate}

\subsection{Error Analysis}
Common prediction errors include:
\begin{itemize}
    \item \textbf{Length Bias}: Both models tend to underpredict scores for concise essays (<100 words)
    \item \textbf{Creative Writing}: Unconventional styles receive lower scores than human raters assign
    \item \textbf{Domain Vocabulary}: Advanced terminology sometimes inflates scores despite weak argumentation
\end{itemize}

\newpage
\section{Conclusion and Future Work}

\subsection{Project Summary}
This project successfully implemented a comprehensive automated essay scoring system achieving research-grade performance on the ASAP dataset. The BERT model's QWK of 0.782 approaches published state-of-the-art results (0.80-0.85), while the multi-dimensional rubric provides interpretable feedback across four writing dimensions.

\subsection{Key Achievements}
\begin{itemize}
    \item Implemented two neural scoring models with comparative evaluation
    \item Developed multi-dimensional rubric scorer (Grammar, Coherence, Relevance, Argument)
    \item Created production-ready Flask web application with REST API
    \item Achieved 77\% error reduction using BERT vs BiLSTM baseline
    \item Integrated GPT-2 feedback generation for constructive guidance
\end{itemize}

\subsection{Limitations}
\begin{itemize}
    \item Single essay set training limits cross-prompt generalization
    \item Length bias toward longer essays
    \item Creative writing styles penalized
    \item Computational requirements (GPU needed for BERT training)
\end{itemize}

\subsection{Future Work}
\begin{enumerate}
    \item \textbf{Cross-Prompt Generalization}: Train on multiple essay sets simultaneously
    \item \textbf{Explainable AI}: Add SHAP analysis for prediction interpretation
    \item \textbf{Bias Mitigation}: Analyze and reduce demographic biases
    \item \textbf{Multilingual Support}: Extend to non-English essays using mBERT
    \item \textbf{Fine-Grained Feedback}: Implement sentence-level scoring
\end{enumerate}

\subsection{Conclusion}
Modern NLP techniques, particularly pretrained transformers like BERT, achieve near-human-level performance on automated essay scoring. The combination of neural scoring with rule-based rubric evaluation provides both accuracy (QWK=0.782) and interpretability, making the system suitable for real-world educational applications.

\newpage
\section*{References}
\addcontentsline{toc}{section}{References}

\begin{enumerate}
    \item Devlin, J., Chang, M. W., Lee, K., \& Toutanova, K. (2019). BERT: Pre-training of Deep Bidirectional Transformers for Language Understanding. \textit{NAACL-HLT 2019}.
    
    \item Taghipour, K., \& Ng, H. T. (2016). A Neural Approach to Automated Essay Scoring. \textit{EMNLP 2016}.
    
    \item The Hewlett Foundation. (2012). Automated Student Assessment Prize (ASAP) Dataset. Kaggle.
    
    \item Vaswani, A., et al. (2017). Attention is All You Need. \textit{NeurIPS 2017}.
    
    \item Hochreiter, S., \& Schmidhuber, J. (1997). Long Short-Term Memory. \textit{Neural Computation, 9(8)}.
\end{enumerate}

\newpage
\section*{AI Usage Declaration}
\addcontentsline{toc}{section}{AI Usage Declaration}

This technical report was prepared by the project team with limited assistance from AI tools for minor formatting and grammar checking purposes only. All technical work, including system design, implementation, experimentation, and analysis, was conducted entirely by the team members.

\textbf{Limited AI Assistance:}
\begin{itemize}
    \item Grammar and spelling correction
    \item LaTeX syntax checking
    \item Minor formatting suggestions
\end{itemize}

\textbf{Original Work by Team:}
\begin{itemize}
    \item Complete system architecture design
    \item All code implementation (preprocessing, models, web app)
    \item Model training and hyperparameter tuning
    \item Experimental design and evaluation
    \item Data analysis and interpretation
    \item Mathematical formulations and explanations
    \item All technical writing and content creation
\end{itemize}

The reported experimental results are authentic and were obtained through our own model training. All code is original and reproducible.

\end{document}
